\section{Introduction}

Automated software-engineering pipelines built on top of Small Language Models (SLMs) are increasingly deployed as multi-stage systems: a model generates a specification (docstring) from code, a second model generates unit tests from the specification, and a third model reconstructs code from the specification and tests. Practitioners assembling such pipelines face a concrete question with no empirically-grounded answer: \emph{should I use the same SLM for documentation, test generation, and code synthesis, or should I specialize by stage?} Existing evaluations of LLM-based software-engineering tools either report on individual stages in isolation (LLM-for-docstring-generation~\citep{Liu2024RAGSurvey}; LLM-for-test-generation~\citep{Schafer2023LLMTest, Tufano2022LLMTest}; LLM-for-code-synthesis~\citep{Chen2021HumanEval}) or evaluate self-consistency within a single LLM across multiple stages~\citep{Allamanis2024RTC, Zhang2025CodingTriangle}. \emph{No published study has measured cross-stage semantic preservation when stages are owned by different LLM families, under a software-engineering-validated defect-detection metric.} The present study addresses this gap.

\subsection{The docstring--test--code closure question}

For each Python function $f$ in a benchmark, three artifacts are canonically associated: the code $C$ (the function body), the docstring $D$ (a natural-language specification), and the reference test suite $T$ (a set of pytest cases exercising~$C$). Together these artifacts form a triangle, and each edge of the triangle can be traversed by a generative pipeline. The pipeline is said to \emph{close} the triangle on a traversal when the reconstructed artifact is semantically equivalent to the original.

We define three closure paths:
\begin{itemize}
  \item \textbf{Path 1 (}$C \to D \to T$\textbf{):} an SLM generates a docstring from the code; a second SLM generates tests from that docstring; we measure whether those tests catch defects introduced in the original code.
  \item \textbf{Path 2 (}$D \to T \to C$\textbf{):} an SLM generates tests from the original docstring; a second SLM reconstructs code from the docstring and generated tests; we measure whether the reconstructed code passes the original reference tests.
  \item \textbf{Path 3 (}$C \to T \to D$\textbf{):} an SLM generates tests from the code; a second SLM recovers a docstring from those tests; we measure whether the recovered docstring is semantically equivalent to the original.
\end{itemize}

Each path is measured by an SE-validated automated metric (mutation kill rate on Path~1, reference-test pass rate on Path~2, docstring BERTScore F1 on Path~3) paired with an external judge SLM's semantic-equivalence rating.

The central research question is whether the mapping of specific SLMs to specific stages of the pipeline matters --- and if so, in what direction, at what magnitude, and with what per-stage attribution.

\subsection{Research questions}

This paper addresses four research questions. A fifth question motivated by the same DOE — comparing open-weight closure rates to a closed-weight frontier ceiling — is outside the scope of the present study and is discussed as a limitation in Section~6 (Threats to Validity) and as future work in Section~7.

\textbf{RQ1.} When the docstring--test--code triangle is traversed by heterogeneous SLMs (different SLM families assigned to the three stages), does the closure rate differ significantly from same-family closure, after controlling for per-sample difficulty?

\textbf{RQ2.} Does mutation kill rate --- measured against the original code's reference-suite mutations --- agree with an external judge SLM's semantic-equivalence rating on each closure path? Where does the two-signal validity policy of Eq.~\ref{eq:valid1}--\ref{eq:valid3} produce disagreement, and does the disagreement structure suggest that the automated metric is capturing the intended signal?

\textbf{RQ3.} Across the $|\mathcal{M}| \times |\mathcal{S}| = 6 \times 3 = 18$ (model, stage) pairings, which stage is the closure bottleneck? Does the bottleneck attribution shift across source benchmark (HumanEval versus MBPP) or across mutation-operator family?

\textbf{RQ4.} What is the false-closure rate --- the proportion of pipelines where the automated metric threshold is satisfied but the judge SLM disagrees --- and is this rate bounded below the empirical noise floor established by the word-shuffled null cell N1?

\subsection{Contributions}

This paper makes five contributions to the empirical literature on multi-SLM software-engineering pipelines:

\begin{enumerate}
  \item \textbf{First empirical study of heterogeneous multi-SLM round-trip closure} across the docstring--test--code triangle, with a pre-registered 20-cell design of experiments spanning mono (6~cells), hetero (11~cells), and null (3~cells) strata. All 20 cells and the aligned concept note were pre-registered in the replication package approximately four weeks before the first sweep result was recorded; the pre-registration is verifiable from the timestamped version history of the released package. Our pre-registration practice follows the emerging guidance of \citet{Sallou2025PreregistrationSE} on pre-registration in empirical software engineering; the open-science replication commitments follow the community roadmap of \citet{Panichella2024OpenScienceSE}.

  \item \textbf{Pre-registered hypothesis validation at per-operator resolution.} The H1 composition (phi-4 spec + qwen-coder test/code) was pre-registered in the design-of-experiments module four weeks before the first sweep result was recorded. The registration includes an explicit per-operator prediction (phi-4 dominates predicate reasoning; qwen-coder dominates comparison operators) drawn from a prior single-stage mutation-testing analysis of LLM-generated tests. H1's empirical per-operator kill rate matches the registered prediction: H1 dominates 4 of 5 mutation operators, with the largest gain ($\Delta = +0.133$) on arithmetic operators. To our knowledge this is the first cross-stage pre-registered LLM composition validated by data at operator resolution.

  \item \textbf{The phi-4-in-test-stage pathology.} We identify a specific stage-model pairing (phi-4 as $L_{\textsf{test}}$) as producing systematically defective test suites at approximately $50\%$ of samples, replicated across three independent cells (M2 mono $= 44\%$, H2 $= 45\%$, H10 $= 55\%$). At operator resolution the pathology localises to the exact mutation families phi-4 mono itself is competitive on --- a cross-stage interaction effect invisible in single-stage evaluations.

  \item \textbf{Judge--metric disagreement decomposition.} On Path~3 (docstring--docstring closure), BERTScore F1 is statistically uncorrelated with the judge SLM's semantic-equivalence rating ($r_3 = 0.02$, $p = 0.359$, $n = 1{,}791$). Stratifying by source benchmark reveals that the disagreement flips direction: on MBPP, BERTScore over-credits equivalence in $29.2\%$ of rows; on HumanEval, BERTScore under-credits in $30.0\%$. This structural directional signature supports the specific claim that BERTScore functions as a surface-form similarity signal rather than a semantic-equivalence signal on our data.

  \item \textbf{Open replication package} including the pre-registered DOE, the full 5{,}890-row sweep TSV, the closure validity decision as a tested Python function (\texttt{closure\_decision.decide\_validity} with 43 unit tests covering all five decision reasons), the test-filter validity gate as a tested Python function (\texttt{closure\_metrics.filter\_tests\_with\_reason} with 8 unit tests covering all four filter reasons), 16 additional wiring-integration tests covering the sweep-pipeline entry points (67 tests total), all G1--G6 gap-closer scripts, 22 LaTeX tables, and 5 figures (1 TikZ framework diagram plus 4 empirical plots). The replication package is released under the MIT license.
\end{enumerate}

A sixth methodological contribution is the demonstration that open-weight SLM evaluation with pre-registered statistical rigor, transparent computational-cost accounting, and open-science reproducibility artifacts~\citep{Panichella2024OpenScienceSE} is a viable methodology for empirical software-engineering research. The six pipeline SLMs used here are all under 30~B parameters and all released with open weights; the entire 5{,}890-row sweep fit within a single-researcher Colab Pro+ monthly allowance ($\sim$530 GPU-hours across A100 and T4, an estimated $25$~kg~CO$_2$-equivalent following~\citet{Lacoste2019MLCO2Emissions} --- see Ethics statement).

\subsection{Roadmap}

The remainder of this paper is organized as follows. Section~2 reviews related work on LLM-based test and docstring generation, round-trip and consistency evaluation, LLM-as-judge protocols, and multi-agent architectural specialization. Section~3 presents the methods, including the notation (Section~3.1), the formal definition of the three closure paths (Section~3.2, Eqs.~\ref{eq:path1}--\ref{eq:path3}), the closure validity decision (Section~3.3, Eq.~\ref{eq:valid1} and Algorithm~\ref{alg:validity}), the test-filter validity gate (Section~3.4, Eq.~\ref{eq:filter} and Algorithm~\ref{alg:test_filter}), the closure metrics (Section~3.5), the 20-cell design of experiments (Section~3.6, Eq.~\ref{eq:doe}), the datasets (Section~3.7), and the statistical methodology (Section~3.8). Section~4 reports the empirical results across the four research questions with anchor tables and figures. Section~5 interprets the findings against RQ1--RQ4, characterizes the phi-4-in-test-stage pathology and the BERTScore surface-form finding, and derives three design guidelines for practitioners deploying multi-SLM SE pipelines. Section~6 details the threats to validity. Section~7 concludes with future research directions.
